\chapter{Method}
\label{ch:method}

\section{Three-layer architecture}
The framework separates real-time execution from high-level reasoning in
three layers (Figure~\ref{fig:architecture}).
\begin{description}
\item[Low-level control and execution layer] The existing deterministic PID
  loops track the active setpoints at a sub-second rate and are never
  generated or modified by the LLM.
\item[Executable heuristic supervisor] A deterministic Python function that
  runs locally at a slower rate, 1--5 minutes in the grinding circuit and
  3--10~s in the toy examples. It monitors a window of recent telemetry and
  updates the setpoints; in the toy examples it also raises one anomaly flag
  per loop.
\item[LLM meta-supervisor] Runs offline. It reads long-term logs of
  performance, disturbances and failure modes and synthesizes new heuristic
  code for the supervisor. It is the only layer that calls the LLM's API.
\end{description}
The LLM's cost, latency and non-determinism therefore never affect the
real-time loop. Before generated code is loaded into the executable
supervisor, it must pass the security check in Section~\ref{sec:sandbox} and
improve on the current supervisor in the evaluation of
Section~\ref{sec:scoring}. Together, the two upper layers take the place of
the TD-MPC supervisory controller.

\begin{figure}[tb]
  \centering
  \includegraphics[width=\textwidth]{three_layer_architecture}
  \caption{The three-layer architecture, shown for the grinding circuit. The
    LLM meta-supervisor works offline on logged data and synthesizes
    heuristic code, which the executable supervisor runs locally without API
    calls. In the toy examples, the process is a tank simulation and the
    supervisor runs every 3--10~s.}
  \label{fig:architecture}
\end{figure}

\section{Supervisor interface}
Every supervisor implements the same function,
\begin{center}
\texttt{supervise(telemetry\_window, active\_setpoints, nominal\_targets)},
\end{center}
where the telemetry window holds, for each controlled tank and each sample,
the level, the output of the loop regulating it (pump effort) and the control
error. The function returns a free-text diagnosis, one adjusted setpoint per
loop (clamped to an allowed range) and one boolean anomaly flag per loop. It
is stateless: it is reloaded for every run and sees only its arguments.
The returned flags are held until the next decision. The setpoint-coordination
experiments in Chapter~\ref{ch:coord} use a variant of this interface
without anomaly flags.

\section{Security sandbox}
\label{sec:sandbox}
Generated code is treated as untrusted and passes three checks before it runs
on any telemetry:
\begin{enumerate}
\item a static check of the abstract syntax tree that rejects imports,
  \texttt{exec}, \texttt{eval}, file and network access, access to dunder
  attributes, oversized sources and anything other than a single top-level
  \texttt{supervise} function with the required signature;
\item execution in a namespace with an explicit whitelist of built-in
  functions and the \texttt{math} and \texttt{statistics} modules, instead of
  Python's real built-ins;
\item a wall-clock timeout on every call (0.5~s), so that an infinite loop
  cannot stall the simulation.
\end{enumerate}
An exception or timeout during evaluation is recorded as a failure, the
setpoints are held and both anomaly flags are raised.

\section{Evaluation and scoring}
\label{sec:scoring}
Each candidate is simulated in closed loop on a fixed \emph{development}
battery of fault scenarios: a fault-free run, single leaks of different onset,
size and duration, simultaneous leaks, persistent severe leaks and, for the
single- and two-tank processes, sensor noise. For a scenario $s$ the score is
\begin{equation}
  J_s = \mathrm{IAE}_s + 500\,V_s + 300\,M_s + c_{\mathrm{FP}}\,F_s
        + 10^4\,E_s + 200\,R_s ,
  \label{eq:score}
\end{equation}
where $\mathrm{IAE}_s$ is the integrated absolute tracking error of all
loops, $V_s$ the number of samples outside the safety bounds, $M_s$ and $F_s$
the number of samples, summed over loops, in which a held anomaly flag
disagrees with the true fault state (missed anomalies and false positives),
$E_s$ the number of exceptions and timeouts, and $R_s$ the distance between
the final and the nominal setpoints for loops whose fault has cleared. The
battery score $J$ is the average of $J_s$ over the scenarios; lower is
better.

Early candidates found several ways to improve $J$ while behaving worse,
which led to three safeguards:
\begin{itemize}
\item \emph{Asymmetric false-positive cost.} A false positive on a loop that
  never faults in the scenario costs $c_{\mathrm{FP}} = 800$ (250 on the
  quadruple-tank process); otherwise $c_{\mathrm{FP}} = 160$. Without this,
  flagging constantly was a cheap way to avoid missed anomalies.
\item \emph{Per-scenario regression guard.} A candidate is promoted only if it
  improves $J$ \emph{and} no single scenario gets worse than the champion by
  more than 5 points (fault-free scenario) or the larger of 100 points and
  15\% (fault scenarios). Averaging otherwise let a candidate win by breaking
  one scenario badly.
\item \emph{Held-out battery.} A second battery with the same scenario types
  but different fault timing and magnitude is never used for promotion. Its
  score is reported to the LLM only as a generalization gap.
\end{itemize}
For the quadruple-tank process an \emph{oracle} is also evaluated. It sets
each flag to the true fault state at every decision and keeps the nominal
setpoints, which is the best detection possible at the given decision
interval. It gives a lower bound that makes it possible to tell whether a
plateau is caused by the search or by the test bed.

\section{Training loop}
Each generation of the LLM meta-supervisor performs the following steps.
\begin{enumerate}
\item Build a prompt. The part that is the same in every call contains a
  description of the process at the level a plant engineer would know it
  (layout, mass balances, parameters, loop pairing, operating point and fault
  types), measured statistics of healthy operation, the calling contract, the
  exact list of allowed built-in functions, the scoring rules, the task and
  the output format. The part that changes contains the champion's source
  code, its per-scenario results, one line per decision for its two
  worst-scoring scenarios (true fault state against returned flags, errors,
  pump efforts and setpoints), a summary of the previous generation's
  rejected candidates, and a list of lessons learned.
\item Sample $K$ candidates from the LLM in parallel. Each response is a JSON
  object with a failure analysis, a numeric self-check of the thresholds
  against the healthy statistics, the new code and at most two new lessons
  of at most 30 words each.
\item Check, load and score every candidate, and promote the best one that
  passes the regression guard.
\end{enumerate}
The quadruple-tank experiments use DeepSeek's \texttt{deepseek-flash} model
through its API, with a configurable reasoning effort (none, low or high),
and the full loop with $K=4$. The loop was built up incrementally: the
single- and two-tank experiments used an earlier version with one candidate
per trial ($K=1$), aggregate failure counts instead of decision traces and
only a short qualitative description of the process, and they were run while
the setup was still changing, with more than one DeepSeek model.
An example of a generated supervisor is given in Appendix~\ref{app:supervisor}.
\TODO{State the prompt size (about 6\,500 tokens on the quadruple-tank process).}
